% Draft response for potential referee comments regarding the absence of DFT+U calculations
\section*{Potential Comment: Why were DFT+U corrections not used for systems containing Cr, Co, Ni, and Fe?}

\textbf{Draft Response:}
We thank the reviewer for raising this important point regarding the treatment of localized $3d$ electrons. We note that while Hubbard corrections (DFT+$U$) are often employed to accurately describe the absolute electronic structure of transition metals, standard PBE has been shown in previous literature to reliably capture the essential structural trends, relative electronic reconstructions, and adsorption energetics in related 2D magnetic materials and functionalized CrCl$_3$ systems. Because the primary focus of this work is to evaluate the \textit{relative} changes in structural, electronic, and catalytic properties induced by chemical functionalization—rather than absolute gap values or correlation effects—standard PBE without Hubbard $U$ corrections was adopted. This level of theory provides a consistent and computationally efficient framework to compare the trends across the different transition-metal dopants.
